\begin{statementbox}
    设 $x^2=2py$（$p>0$）上两个不同切点为 $A,B$，两切线交于 $P$；
    $M$ 为弦 $AB$ 中点。核心关系如下图。

    \begin{center}
\begin{tikzpicture}[x=1.05cm,y=1.05cm,font=\small,>=Stealth]
\path (-2,-2.1) rectangle (9,3.6);
\coordinate (A) at (-1,.25); \coordinate (B) at (3,2.25);
\coordinate (P) at (1,-.75); \coordinate (M) at (1,1.25);
\coordinate (N) at (1,.25); \coordinate (H) at (.2,.85);
\fill[blue!12] (P)--(A)--(M)--cycle;
\fill[orange!20] (P)--(M)--(B)--cycle;
\draw[->,gray] (-1.8,0)--(3.8,0) node[right] {$x$};
\draw[->,gray] (0,-1.35)--(0,3.2) node[above] {$y$};
\draw[blue!70!black,thick,domain=-1.6:3.45,samples=100] plot (\x,{\x*\x/4});
\draw[gray!80,domain=-1.7:1.5] plot (\x,{-.5*\x-.25});
\draw[gray!80,domain=.7:3.4] plot (\x,{1.5*\x-2.25});
\draw[very thick] (A)--(B);
\draw[thick] (A)--(P)--(B);
\draw[orange!80!black,very thick] (P)--(M);
\foreach \pt/\pos in {B/above left,M/above left}
  \fill (\pt) circle (1.6pt) node[\pos] {$\pt$};
\fill (A) circle (1.6pt); \node[anchor=east] at (-1.3,-.35) {$A$};
\fill (P) circle (1.6pt); \node[anchor=west] at (1.4,-.65) {$P$};

\fill[blue!70!black] (N) circle (1.6pt) node[right] {$N$};
\draw[dashed] (A)--(-1,-1.25) (B)--(3,-1.25);
\draw[<->] (-1,-1.2)--(3,-1.2) node[midway,below] {$|x_2-x_1|$};

\node[anchor=north west,align=left,text width=4.7cm] at (4.1,3.1) {\textbf{两切线交点与面积}\\[6pt]
$x_P=x_M=\dfrac{x_1+x_2}{2}$\\[8pt]
$PM=\dfrac{(x_1-x_2)^2}{4p}$\\[8pt]
$S=\dfrac12 PM\,|x_2-x_1|$\\[8pt]
蓝、橙两块拼成 $\triangle PAB$。\\[6pt]
$PN=NM$；$N$ 在曲线上。\\[6pt]
图取 $p=2$，\\[3pt]
$x_1=-1,\ x_2=3$。};
\end{tikzpicture}
\end{center}


    \textbf{\textcolor{CStatement}{条件}}
    \begin{description}[style=nextline, leftmargin=2.8em, labelsep=0.5em, itemsep=0.4em]
        \item[\textbf{\textcolor{CStatement}{条件一（曲线与切点）：}}]
        设抛物线为 $x^2=2py\;(p>0)$，取其上两个\textbf{不同}的点
        \[
            A(x_1,y_1),\qquad B(x_2,y_2),\qquad
            y_1=\frac{x_1^2}{2p},\quad y_2=\frac{x_2^2}{2p}.
        \]
        分别在 $A,B$ 处作切线，两切线交于 $P(x_0,y_0)$。
        因为抛物线上的点由横坐标唯一确定，所以 $A\ne B$ 等价于 $x_1\ne x_2$。
        \item[\textbf{\textcolor{CStatement}{条件二（三角形）：}}]
        以切点弦 $AB$ 为底边、两切线交点 $P$ 为顶点，构成三角形 $PAB$。
        本结论沿用项目中的名称“抛物线阿基米德三角形”。
    \end{description}

    \textbf{\textcolor{CStatement}{四个主结论}}
    \begin{description}[style=nextline, leftmargin=2.8em, labelsep=0.5em, itemsep=0.55em]
        \item[\textbf{\textcolor{CStatement}{结论一（两切线交点）：}}]
        \[
            \boxed{P\left(\frac{x_1+x_2}{2},\,\frac{x_1x_2}{2p}\right).}
        \]
        横坐标是两切点横坐标的平均数；纵坐标与其乘积有关。

        \item[\textbf{\textcolor{CStatement}{结论二（中线与长度）：}}]
        设 $M$ 是 $AB$ 的中点，则
        \[
            M\left(\frac{x_1+x_2}{2},\,\frac{x_1^2+x_2^2}{4p}\right),
            \qquad \boxed{PM=\frac{(x_1-x_2)^2}{4p}.}
        \]
        $P$ 与 $M$ 横坐标相同，所以 $PM$ 是\textbf{竖直线段}，
        与抛物线的对称轴（$y$ 轴）平行；且 $M$ 位于 $P$ 上方。

        \item[\textbf{\textcolor{CStatement}{结论三（面积）：}}]
        \[
            \boxed{S_{\triangle PAB}=\frac{|x_1-x_2|^3}{8p}.}
        \]
        几何理解：\textbf{面积 = 竖直中线 $PM$ 的一半 $\times$ 底边的水平投影长度}，即
        \[
            S_{\triangle PAB}=\frac12\,PM\cdot |x_1-x_2|.
        \]
        注意：此处\textbf{不是}把 $PM$ 当作 $AB$ 上的高。

        \item[\textbf{\textcolor{CStatement}{结论四（切点弦）：}}]
        若 $P(x_0,y_0)$，则
        \[
            \boxed{AB:\quad x_0x=p(y+y_0).}
        \]
        由此可反求切点横坐标，它们是方程
        \[
            x^2-2x_0x+2py_0=0
        \]
        的两个实数解。存在两个\textbf{不同}切点的条件为
        \[
            \boxed{x_0^2-2py_0>0.}
        \]
    \end{description}

    \textbf{\textcolor{CStatement}{一个附加的几何性质}}

    过 $P$ 作平行于 $y$ 轴的直线，与抛物线交于 $N$，则
    \[
        \boxed{PN=NM=\frac{(x_1-x_2)^2}{8p}.}
    \]
    换句话说：\textbf{抛物线在竖直中线上的交点 $N$，恰好是线段 $PM$ 的中点。}


    \textbf{\textcolor{CStatement}{与初中知识的连接}}

    在普通三角形中，\textbf{中线不一定是高}：如果一条中线恰好竖直，
    可以把它当作两个小三角形共同的底边，以水平距离作高。
    本结论先证明 $PM$ 竖直，再用同样的初中面积方法，
    不需要先求斜边 $AB$ 的长度。
    详细的初中例子及与圆的两切线三角形的比较，见“理解”和“例题”。

    \medskip

    \textbf{\textcolor{CStatement}{适用提醒}}

    所有坐标和系数均按 $x^2=2py\;(p>0)$ 写出。对于 $y^2=2px$，应交换
    $x,y$ 的位置重新使用相应公式，不得直接照抄。
\end{statementbox}
